\input{../tex-inputs/latex-leseansicht-vorspann}

\section[Arthur Schnitzler an Theodor von Sosnosky, {[}zwischen 26. 1. 1901 und 29. 1. 1901?{]}, Fragment]{L04237 Arthur Schnitzler an Theodor von Sosnosky, [zwischen 26. 1. 1901 und
               29. 1. 1901?], Fragment}
\nopagebreak\mylabel{L04237v}
\rehead{ }\normalsize\beginnumbering\briefempfaengerindex{Sosnosky, Theodor von@\textsc{Sosnosky, Theodor von}!zzzSchnitzler, Arthur@\emph{von Arthur Schnitzler}!1901-01-292@{{[}zwischen 26. 1. 1901 und 29. 1. 1901?{]}}|(be}
\toendnotes[C]{\smallbreak\pagebreak[2]}
\correspDesc{Versand  durch Arthur Schnitzler im Zeitraum [zwischen 26. 1. 1901 und
                  29. 1. 1901?] in Wien
\newline{}Erhalt  durch Theodor von Sosnosky in Kremsmünster}\toendnotes[C]{\smallbreak}
\buchAlsQuelle{Theodor v. Sosnosky: \emph{Artur Schnitzler als Lyriker.} In: \emph{Neues Wiener Journal}, Jg. 40, Nr. 13.823, 18. 5. 1932, S. 6–7.}\toendnotes[C]{\smallbreak}
\pstart
           \noindent{}Insbeſondere was Sie \label{K_L04237-1v}\edtext{über den »Reigen\pwindex{Schnitzler, Arthur 15.\,5.\,1862 Wien – 21.\,10.\,1931 ebd.@\textsc{Schnitzler, Arthur} (15.\,5.\,1862 Wien – 21.\,10.\,1931 ebd.), \emph{Schriftsteller, Mediziner}!Reigen. Zehn Dialoge@\strich\emph{Reigen. Zehn Dialoge}|pw}« ſagen\pwindex{Sosnosky, Theodor von 4.\,1.\,1866 Budapest – 3.\,2.\,1943 Wien@\textsc{Sosnosky, Theodor von} (4.\,1.\,1866 Budapest – 3.\,2.\,1943 Wien), \emph{Schriftsteller}!Jung-Wien«. Studienblätter@\strich\emph{»Jung-Wien«. Studienblätter}|pwv}}{\lemma{\textnormal{\emph{über den »Reigen« sagen}}}\Cendnote{\textnormal{Die Datierung ist im Druck des Briefes
                  nicht überliefert. Schnitzler nimmt Bezug
                  auf Sosnoskys\pwindex{Sosnosky, Theodor von 4.\,1.\,1866 Budapest – 3.\,2.\,1943 Wien@\textsc{Sosnosky, Theodor von} (4.\,1.\,1866 Budapest – 3.\,2.\,1943 Wien), \emph{Schriftsteller}|pwk} Essay \emph{Jung-Wien}\pwindex{Sosnosky, Theodor von 4.\,1.\,1866 Budapest – 3.\,2.\,1943 Wien@\textsc{Sosnosky, Theodor von} (4.\,1.\,1866 Budapest – 3.\,2.\,1943 Wien), \emph{Schriftsteller}!Jung-Wien«. Studienblätter@\strich\emph{»Jung-Wien«. Studienblätter}|pwk}, der in drei Folgen vom 20. 1. bis
                  zum 3. 2. 1901 in den Beilagen der \emph{Norddeutschen Allgemeinen Zeitung}\pwindex{Norddeutsche Allgemeine Zeitung@\emph{Norddeutsche Allgemeine Zeitung}|pwk} erschien (\emph{»Jung-Wien«. Studienblätter von Theodor v.
                        Sosnosky}\pwindex{Sosnosky, Theodor von 4.\,1.\,1866 Budapest – 3.\,2.\,1943 Wien@\textsc{Sosnosky, Theodor von} (4.\,1.\,1866 Budapest – 3.\,2.\,1943 Wien), \emph{Schriftsteller}!Jung-Wien«. Studienblätter@\strich\emph{»Jung-Wien«. Studienblätter}|pwk}. In: Beilage zur \emph{Norddeutschen Allgemeinen Zeitung}\pwindex{Norddeutsche Allgemeine Zeitung@\emph{Norddeutsche Allgemeine Zeitung}|pwk}, Jg. 40, Nr. 17a,
                        20. 1. 1901, S. [1]–2, Nr. 20, 26. 1. 1901, S. [1]
                     und Nr. 29a, 3. 2. 1901, S. [1]–2). Um den \emph{Reigen}\pwindex{Schnitzler, Arthur 15.\,5.\,1862 Wien – 21.\,10.\,1931 ebd.@\textsc{Schnitzler, Arthur} (15.\,5.\,1862 Wien – 21.\,10.\,1931 ebd.), \emph{Schriftsteller, Mediziner}!Reigen. Zehn Dialoge@\strich\emph{Reigen. Zehn Dialoge}|pwk} geht es im zweiten Teil, der am
                     26. 1. 1901 publiziert wurde. Der Brief wurde also frühestens
                     an diesem Tag abgefasst und – mit Blick auf die Postlaufzeit –
                  spätestens am 29. 1. 1901, denn am 30. 1. 1901
                  antwortete Sosnosky\pwindex{Sosnosky, Theodor von 4.\,1.\,1866 Budapest – 3.\,2.\,1943 Wien@\textsc{Sosnosky, Theodor von} (4.\,1.\,1866 Budapest – 3.\,2.\,1943 Wien), \emph{Schriftsteller}|pwk} aus Kremsmünster\oindex{Kremsmünster@\textbf{Kremsmünster}, \emph{Verwaltungsgebiet}|pwk} darauf.}}}\label{K_L04237-1} – von Ihrer liebenswürdigen,
               meiner eigenen Empfindung nach nicht{ }ſelten zu liebenswürdigen Beurteilung mancher
                  \label{K_L04237-2v}\edtext{meiner anderen Schriften}{\lemma{\textnormal{\emph{meiner anderen Schriften}}}\Cendnote{\textnormal{Außer \emph{Reigen}\pwindex{Schnitzler, Arthur 15.\,5.\,1862 Wien – 21.\,10.\,1931 ebd.@\textsc{Schnitzler, Arthur} (15.\,5.\,1862 Wien – 21.\,10.\,1931 ebd.), \emph{Schriftsteller, Mediziner}!Reigen. Zehn Dialoge@\strich\emph{Reigen. Zehn Dialoge}|pwk} nennt Sosnosky\pwindex{Sosnosky, Theodor von 4.\,1.\,1866 Budapest – 3.\,2.\,1943 Wien@\textsc{Sosnosky, Theodor von} (4.\,1.\,1866 Budapest – 3.\,2.\,1943 Wien), \emph{Schriftsteller}|pwk}{ }\emph{Anatol}\pwindex{Schnitzler, Arthur 15.\,5.\,1862 Wien – 21.\,10.\,1931 ebd.@\textsc{Schnitzler, Arthur} (15.\,5.\,1862 Wien – 21.\,10.\,1931 ebd.), \emph{Schriftsteller, Mediziner}!Anatol@\strich\emph{Anatol}|pwk}, \emph{Liebelei}\pwindex{Schnitzler, Arthur 15.\,5.\,1862 Wien – 21.\,10.\,1931 ebd.@\textsc{Schnitzler, Arthur} (15.\,5.\,1862 Wien – 21.\,10.\,1931 ebd.), \emph{Schriftsteller, Mediziner}!Liebelei. Schauspiel in drei Akten@\strich\emph{Liebelei. Schauspiel in drei Akten}|pwk}, \emph{Drei
                     Elxire}\pwindex{Schnitzler, Arthur 15.\,5.\,1862 Wien – 21.\,10.\,1931 ebd.@\textsc{Schnitzler, Arthur} (15.\,5.\,1862 Wien – 21.\,10.\,1931 ebd.), \emph{Schriftsteller, Mediziner}!drei Elixire@\strich\emph{Die drei Elixire}|pwk}, \emph{Sterben}\pwindex{Schnitzler, Arthur 15.\,5.\,1862 Wien – 21.\,10.\,1931 ebd.@\textsc{Schnitzler, Arthur} (15.\,5.\,1862 Wien – 21.\,10.\,1931 ebd.), \emph{Schriftsteller, Mediziner}!Sterben. Novelle@\strich\emph{Sterben. Novelle}|pwk}, \emph{Ein Abschied}\pwindex{Schnitzler, Arthur 15.\,5.\,1862 Wien – 21.\,10.\,1931 ebd.@\textsc{Schnitzler, Arthur} (15.\,5.\,1862 Wien – 21.\,10.\,1931 ebd.), \emph{Schriftsteller, Mediziner}!Abschied@\strich\emph{Ein Abschied}|pwk} und \emph{Die Toten
                     schweigen}\pwindex{Schnitzler, Arthur 15.\,5.\,1862 Wien – 21.\,10.\,1931 ebd.@\textsc{Schnitzler, Arthur} (15.\,5.\,1862 Wien – 21.\,10.\,1931 ebd.), \emph{Schriftsteller, Mediziner}!Toten schweigen@\strich\emph{Die Toten schweigen}|pwk}.}}}\label{K_L04237-2} ganz zu{ }ſchweigen –, hat mich ganz außerordentlich
               gefreut {\dots} Die Keuſchheit des {\dotsfour}Blattes\pwindex{Neues Wiener Tagblatt@\emph{Neues Wiener Tagblatt}|pwv} (das eine Beſprechung
               des Buches\pwindex{Schnitzler, Arthur 15.\,5.\,1862 Wien – 21.\,10.\,1931 ebd.@\textsc{Schnitzler, Arthur} (15.\,5.\,1862 Wien – 21.\,10.\,1931 ebd.), \emph{Schriftsteller, Mediziner}!Reigen. Zehn Dialoge@\strich\emph{Reigen. Zehn Dialoge}|pwv}, die ich ihm
               angeboten, trotz{ }ſeines zur Schau getragenen »Freiſinns« aus Gründen der Sittlichkeit
               abgelehnt hatte) hat mich ebenſo amüſiert wie Sie.\pend
           \selectlanguage{ngerman}\endnumbering\briefempfaengerindex{Sosnosky, Theodor von@\textsc{Sosnosky, Theodor von}!zzzSchnitzler, Arthur@\emph{von Arthur Schnitzler}!1901-01-262@{{[}zwischen 26. 1. 1901 und 29. 1. 1901?{]}}|)be}\mylabel{L04237h}
\begin{anhang}
\end{anhang}\newcommand{\dateiname}{L04237}\newcommand{\titel}{Arthur Schnitzler an Theodor von Sosnosky, [zwischen 26. 1. 1901 und 29. 1. 1901?], Fragment}\newcommand{\editorInnen}{Herausgegeben von Jahnke, SelmaMüller, Martin Anton}\input{../tex-inputs/latex-leseansicht-abspann}
